\documentclass[11pt]{article}

\usepackage[margin=1in]{geometry}
\usepackage{amsmath,amssymb,amsfonts,bm}
\usepackage{booktabs}
\usepackage{enumitem}
\usepackage{hyperref}

\title{Amortized Alpha Training for the 3D Hexner-Mod Low-Running-Cost Velocity Case}
\author{}
\date{}

\newcommand{\R}{\mathbb{R}}
\newcommand{\E}{\mathbb{E}}
\newcommand{\DeltaI}{\Delta(\mathcal I)}
\newcommand{\softmax}{\operatorname{softmax}}

\begin{document}
\setcounter{MaxMatrixCols}{20}
\maketitle

\section{Purpose and Run Identifier}

This note documents the 3D amortized Hexner-mod experiment stored in
\[
\texttt{runs/hexner\_mod\_3d\_amortized\_vel\_low\_running\_cost}.
\]
The run uses the original target-based 3D Hexner-mod game
(\texttt{HexnerMod3DOriginalGame}) with two hidden payoff types, zero terminal
velocity penalties, and reduced running-control penalties.  The learned object
is an amortized public signaling policy
\[
    f_\psi : (x_0,p_0) \mapsto \alpha,
\]
where $x_0$ is the initial joint state, $p_0$ is the initial public belief over
types, and $\alpha$ is a full public-tree signaling policy.

The checkpoint metadata for this run reports:
\begin{center}
\begin{tabular}{ll}
\toprule
Quantity & Value \\
\midrule
Horizon & $T=1.0$ \\
Number of time steps & $K=10$ \\
Step size & $\tau = T/K = 0.1$ \\
Number of types & $I=2$ \\
Training iterations & $2000$ \\
Final checkpoint step & $1999$ \\
Batch size & $32$ \\
Optimizer & Adam \\
Learning rate & $10^{-3}$ \\
Network hidden layer & one hidden layer of width $64$ \\
Activation & SiLU \\
Dropout & $0$ \\
Training device/dtype & CPU, \texttt{float32} \\
\bottomrule
\end{tabular}
\end{center}

\section{3D Game Setting}

\subsection{State, Controls, and Dynamics}

Each player has a 3D double-integrator state.  For player $j \in \{1,2\}$,
\[
    x_j =
    \begin{bmatrix}
        p_{j,x} & p_{j,y} & p_{j,z} &
        v_{j,x} & v_{j,y} & v_{j,z}
    \end{bmatrix}^{\!\top}
    \in \R^6.
\]
The full state is the stacked vector
\[
    x =
    \begin{bmatrix}
        x_1 \\ x_2
    \end{bmatrix}
    \in \R^{12}.
\]
Player 1's acceleration control is $u \in \R^3$, and player 2's acceleration
control is $v \in \R^3$.

The discrete-time dynamics are
\[
    x_{k+1} = A x_k + B_1 u_k + B_2 v_k,
\]
where
\[
    A =
    \begin{bmatrix}
        A_p & 0 \\
        0 & A_p
    \end{bmatrix},
    \qquad
    B_1 =
    \begin{bmatrix}
        B_p \\ 0
    \end{bmatrix},
    \qquad
    B_2 =
    \begin{bmatrix}
        0 \\ B_p
    \end{bmatrix}.
\]
The per-player double-integrator matrices are
\[
    A_p =
    \begin{bmatrix}
        I_3 & \tau I_3 \\
        0 & I_3
    \end{bmatrix},
    \qquad
    B_p =
    \begin{bmatrix}
        \frac{1}{2}\tau^2 I_3 \\
        \tau I_3
    \end{bmatrix}.
\]
Thus player 1's state is driven only by $u_k$, and player 2's state is driven
only by $v_k$.

\subsection{Hidden Types and Type-Dependent Targets}

There are two hidden payoff types,
\[
    \theta_0=-1,\qquad \theta_1=+1.
\]
The base target vector is
\[
    z =
    \begin{bmatrix}
        1 & 1 & 0 & 0 & 0 & 0
    \end{bmatrix}^{\!\top}.
\]
Therefore the type-dependent target state for each player is
\[
    z_i = \theta_i z.
\]
Explicitly,
\[
    z_0 =
    \begin{bmatrix}
        -1 & -1 & 0 & 0 & 0 & 0
    \end{bmatrix}^{\!\top},
    \qquad
    z_1 =
    \begin{bmatrix}
        1 & 1 & 0 & 0 & 0 & 0
    \end{bmatrix}^{\!\top}.
\]
The $z$-coordinate target is zero for both types, so the landing plane is
$p_z=0$.  The terminal velocity target is also zero in all three velocity
directions.

\subsection{Running Cost and the ``Low Running Cost'' Setting}

The game is zero-sum.  For a fixed type $i$, the one-step running cost has the
form
\[
    \ell_i(u_k,v_k)
    =
    \frac{1}{2} u_k^\top R_i u_k
    -
    \frac{1}{2} v_k^\top S_i v_k.
\]
The matrices are type-independent in this run, so $R_i=R$ and $S_i=S$ for
both types.  The base matrices in code are
\[
    R_{\text{base}} =
    \operatorname{diag}(0.05,\,0.025,\,0.05),
    \qquad
    S_{\text{base}} =
    \operatorname{diag}(0.05,\,0.10,\,0.05).
\]
The run uses
\[
    \texttt{R1\_scale}=0.25,
    \qquad
    \texttt{R2\_scale}=0.10.
\]
The implementation stores the matrices used by the solver as
\[
    R = 2 \cdot \texttt{R1\_scale} \cdot R_{\text{base}},
    \qquad
    S = 2 \cdot \texttt{R2\_scale} \cdot S_{\text{base}}.
\]
Therefore the actual running-cost matrices are
\[
    R =
    \operatorname{diag}(0.025,\,0.0125,\,0.025),
    \qquad
    S =
    \operatorname{diag}(0.010,\,0.020,\,0.010).
\]
Relative to the default original 3D setting, these are lower effort penalties.
The practical effect is that both players can use larger accelerations at a
lower running-cost price.  This is why the run name contains
\texttt{low\_running\_cost}.

The action bounds stored in the config are
\[
    u_k \in [-14,14]^3,
    \qquad
    v_k \in [-14,14]^3.
\]
The current Riccati training path computes the unconstrained LQ saddle policy;
these bounds are available for downstream clipping during rollout if desired.

\subsection{Terminal Cost, Landing Plane, and Zero Final Velocity}

For type $i$, the terminal cost is
\[
    g_i(x_T)
    =
    (x_{1,T}-z_i)^\top K_{1,i}(x_{1,T}-z_i)
    -
    (x_{2,T}-z_i)^\top K_{2,i}(x_{2,T}-z_i).
\]
This says that player 1 is penalized for missing the type-dependent target,
while player 2 enters with the opposite sign because this is a zero-sum game.
For this run,
\[
    K_{1,i}=K_{2,i}=D_i.
\]
The terminal diagonal weights are
\[
    D_0 =
    \operatorname{diag}(1,\,20,\,20,\,20,\,20,\,20),
\]
\[
    D_1 =
    \operatorname{diag}(20,\,1,\,20,\,20,\,20,\,20).
\]
The first three entries correspond to terminal position error
$(p_x,p_y,p_z)$, and the last three entries correspond to terminal velocity
error $(v_x,v_y,v_z)$.

The important velocity-condition setting is
\[
    \texttt{terminal\_velocity\_scale}=20.
\]
Thus every final velocity component has terminal weight $20$.  In expanded
form, for a generic single-player terminal state
\[
    \xi =
    \begin{bmatrix}
        p_x & p_y & p_z & v_x & v_y & v_z
    \end{bmatrix}^{\!\top},
\]
the type-0 single-player target cost is
\[
\begin{aligned}
    \phi_0(\xi)
    &=
    (\xi-z_0)^\top D_0(\xi-z_0) \\
    &=
    (p_x+1)^2
    +20(p_y+1)^2
    +20p_z^2
    +20v_x^2
    +20v_y^2
    +20v_z^2,
\end{aligned}
\]
and the type-1 single-player target cost is
\[
\begin{aligned}
    \phi_1(\xi)
    &=
    (\xi-z_1)^\top D_1(\xi-z_1) \\
    &=
    20(p_x-1)^2
    +(p_y-1)^2
    +20p_z^2
    +20v_x^2
    +20v_y^2
    +20v_z^2.
\end{aligned}
\]
Therefore the full terminal cost is
\[
    g_i(x_T)=\phi_i(x_{1,T})-\phi_i(x_{2,T}).
\]

This terminal objective has three notable features:
\begin{enumerate}[label=(\roman*)]
    \item The target location depends on the hidden type through
    $z_i=\theta_i z$.
    \item The landing plane is type-independent and equals $p_z=0$.
    \item Zero terminal velocity is enforced softly through the large velocity
    weights $20v_x^2+20v_y^2+20v_z^2$.
\end{enumerate}

\subsection{Default Initial State}

The default initial state used in the run metadata is
\[
    x_{\mathrm{ref}} =
    \begin{bmatrix}
    -1 & 0 & 1 & 0 & 0 & 0 &
     1 & 0 & 1 & 0 & 0 & 0
    \end{bmatrix}^{\!\top}.
\]
Thus player 1 starts at position $(-1,0,1)$ with zero velocity, while player 2
starts at $(1,0,1)$ with zero velocity.  Both begin at altitude $p_z=1$ and are
asked to land at $p_z=0$ with zero final velocity.

\section{Public Tree, Signaling Policy, and Belief Updates}

\subsection{Public Prototype Tree}

The game has $I=2$ hidden types and also uses $I=2$ public prototype branches
at each node.  The public tree has depth $K=10$.  A node at depth $k$ is
identified by a public prototype history
\[
    \omega_k=(a_0,a_1,\ldots,a_{k-1}),
    \qquad
    a_t \in \{0,1\}.
\]
There are $2^k$ nodes at depth $k$.  The tree indexer uses the recursion
\[
    \operatorname{child}(k,\omega,a) = 2\omega + a
\]
in local integer indexing.

\subsection{Alpha as a Type-Conditioned Public Signaling Strategy}

At each depth $k$, node $\omega$, and type $i$, the policy chooses a public
prototype $a$ according to
\[
    \alpha_{k,\omega,i}^{a}
    =
    \Pr(a_k=a \mid \omega_k=\omega,\; i).
\]
For each fixed $(k,\omega,i)$,
\[
    \alpha_{k,\omega,i}^{a}\ge 0,
    \qquad
    \sum_{a=0}^{I-1}\alpha_{k,\omega,i}^{a}=1.
\]
Intuitively, $\alpha$ is the signaling strategy of the informed player.  Since
the public branch $a$ is observed, it updates the public belief about the
hidden type.  The subsequent Riccati solve computes the continuous controls
associated with every edge $(k,\omega,a)$.

\subsection{Belief Propagation}

Let $p_{k,\omega}(i)$ denote the public belief at node $(k,\omega)$.  Given
$\alpha_{k,\omega,i}^{a}$, the probability of observing public prototype $a$
at node $(k,\omega)$ is
\[
    \lambda_{k,\omega}^{a}
    =
    \sum_{i=0}^{I-1}
    p_{k,\omega}(i)\alpha_{k,\omega,i}^{a}.
\]
After observing $a$, the posterior belief at the child node $(k+1,\omega a)$
is obtained by Bayes' rule:
\[
    p_{k+1,\omega a}(i)
    =
    \frac{
        p_{k,\omega}(i)\alpha_{k,\omega,i}^{a}
    }{
        \sum_{j=0}^{I-1}p_{k,\omega}(j)\alpha_{k,\omega,j}^{a}
    }.
\]
The implementation also tracks node masses, but the conditional beliefs above
are the quantities used to compute belief-averaged costs and policies.

\section{Inner Tree LQ Game for a Fixed \texorpdfstring{$\alpha$}{alpha}}

For a fixed signaling policy $\alpha$, the belief tree is fixed.  At every
edge $(k,\omega,a)$, the game uses the posterior belief at the child node,
$p_{k+1,\omega a}$, to compute belief-averaged running costs
\[
    \bar R_{k,\omega}^{a}
    =
    \sum_i p_{k+1,\omega a}(i) R_i,
    \qquad
    \bar S_{k,\omega}^{a}
    =
    \sum_i p_{k+1,\omega a}(i) S_i.
\]
At terminal leaves, the type-dependent terminal quadratics are also averaged
under the leaf belief.

The Riccati recursion stores a quadratic value at every node:
\[
    V_{k,\omega}(x)
    =
    \frac{1}{2}x^\top P_{k,\omega}x
    +r_{k,\omega}^\top x
    +c_{k,\omega}.
\]
At an edge $(k,\omega,a)$, the local saddle problem is
\[
\begin{aligned}
    \min_u \max_v\quad
    &\tau\left(
        \frac{1}{2}u^\top \bar R_{k,\omega}^{a} u
        -
        \frac{1}{2}v^\top \bar S_{k,\omega}^{a} v
    \right) \\
    &\quad
    + V_{k+1,\omega a}
    \left(Ax+B_1u+B_2v\right).
\end{aligned}
\]
Solving this LQ saddle problem gives affine feedback laws
\[
    u_{k,\omega}^{a}(x)
    =
    K_{u,k,\omega}^{a}x+\kappa_{u,k,\omega}^{a},
    \qquad
    v_{k,\omega}^{a}(x)
    =
    K_{v,k,\omega}^{a}x+\kappa_{v,k,\omega}^{a}.
\]
The optimized edge value is another quadratic,
\[
    V_{k,\omega}^{a,\mathrm{edge}}(x)
    =
    \frac{1}{2}x^\top P_{k,\omega}^{a}x
    +(r_{k,\omega}^{a})^\top x
    +c_{k,\omega}^{a}.
\]
The node value averages over public branches using the edge probabilities
$\lambda_{k,\omega}^{a}$:
\[
    V_{k,\omega}(x)
    =
    \sum_{a=0}^{I-1}
    \lambda_{k,\omega}^{a}
    V_{k,\omega}^{a,\mathrm{edge}}(x).
\]
This gives the root value
\[
    J(\alpha;x_0,p_0)=V_{0,\emptyset}(x_0).
\]

\section{Amortized Training}

\subsection{From Direct Alpha Optimization to Amortization}

In a non-amortized solve, one would optimize a separate $\alpha$ tensor for
one fixed initial condition $(x_0,p_0)$.  Amortized training instead learns a
single neural network that maps many possible contexts to their corresponding
signaling policies:
\[
    \alpha = f_\psi(x_0,p_0).
\]
The network parameters $\psi$ are shared across all sampled initial states and
beliefs.  At deployment time, a new $(x_0,p_0)$ can be fed through the network
once to produce a full tree policy.

\subsection{Network Parameterization}

The input is the concatenated vector
\[
    h_0 =
    \begin{bmatrix}
        x_0 \\ p_0
    \end{bmatrix}
    \in \R^{12+2}=\R^{14}.
\]
For this run, the network is a one-hidden-layer MLP:
\[
    h_1 = \operatorname{SiLU}(W_1h_0+b_1),
    \qquad
    \ell = W_2h_1+b_2.
\]
The hidden width is $64$.  The output logits $\ell$ are reshaped into
\[
    \ell
    \in
    \R^{K \times N_{\max}\times I\times I},
    \qquad
    K=10,\quad I=2,\quad N_{\max}=2^{10}=1024.
\]
Only the first $2^k$ node entries are active at depth $k$; the extra entries
exist only so that the tensor has a rectangular shape.  The signaling
probabilities are obtained by a softmax over the final action/prototype axis:
\[
    \alpha_{k,\omega,i}^{a}
    =
    \frac{
        \exp(\ell_{k,\omega,i}^{a})
    }{
        \sum_{b=0}^{I-1}\exp(\ell_{k,\omega,i}^{b})
    }.
\]

\subsection{Sampling Initial States}

Each training batch contains $32$ contexts.  The reference state is
\[
    x_{\mathrm{ref}} =
    \begin{bmatrix}
    -1 & 0 & 1 & 0 & 0 & 0 &
     1 & 0 & 1 & 0 & 0 & 0
    \end{bmatrix}^{\!\top}.
\]
For most samples, the position components are perturbed uniformly:
\[
    p_{1,0} = p_{1,\mathrm{ref}} + \delta_1,
    \qquad
    \delta_1 \sim \operatorname{Unif}([-0.4,0.4]^3),
\]
\[
    p_{2,0} = p_{2,\mathrm{ref}} + \delta_2,
    \qquad
    \delta_2 \sim \operatorname{Unif}([-0.4,0.4]^3).
\]
The run uses
\[
    \texttt{p1\_pos\_jitter}=0.4,
    \qquad
    \texttt{p2\_pos\_jitter}=0.4.
\]
Velocity jitter is disabled:
\[
    \texttt{p1\_vel\_jitter}=0,
    \qquad
    \texttt{p2\_vel\_jitter}=0.
\]
Thus the initial velocities remain zero in the sampled training contexts.

\subsection{Sampling Initial Beliefs}

For $I=2$, the prior belief is sampled by drawing the probability of type $0$:
\[
    p_0(0) \sim \operatorname{Unif}(0.05,0.95),
    \qquad
    p_0(1)=1-p_0(0).
\]
The lower bound $0.05$ is set by
\[
    \texttt{prior\_min}=0.05.
\]
This avoids extremely degenerate priors and trains the network over a broad
range of uncertainty levels.

Additionally, with probability
\[
    \texttt{default\_context\_prob}=0.15,
\]
a sample in the batch is reset to the default reference context
\[
    (x_{\mathrm{ref}},\, [0.5,0.5]).
\]
This keeps the canonical scenario represented throughout training while still
learning over many nearby initial states and beliefs.

\subsection{Training Objective}

For each sampled context $(x_0^{(b)},p_0^{(b)})$, the network predicts
\[
    \alpha^{(b)} = f_\psi(x_0^{(b)},p_0^{(b)}).
\]
The belief tree is built from $(\alpha^{(b)},p_0^{(b)})$, the belief-averaged
costs are computed, and the tree-structured Riccati recursion evaluates the
root value
\[
    J_\psi^{(b)}
    =
    J(\alpha^{(b)};x_0^{(b)},p_0^{(b)})
    =
    V_{0,\emptyset}^{(b)}(x_0^{(b)}).
\]
The mini-batch objective is the mean root value:
\[
    \mathcal L(\psi)
    =
    \frac{1}{B}
    \sum_{b=1}^{B}
    J(f_\psi(x_0^{(b)},p_0^{(b)});x_0^{(b)},p_0^{(b)}),
    \qquad B=32.
\]
Training performs gradient descent on $\mathcal L(\psi)$ using Adam with
learning rate $10^{-3}$ and no weight decay.  Gradients flow through the
network, belief construction, averaged-cost computation, and Riccati root-value
computation.

\section{Interpretation of the Low-Running-Cost Velocity Run}

The key choices in this run are:
\begin{itemize}
    \item Reduced running-control penalties:
    \[
        R=\operatorname{diag}(0.025,0.0125,0.025),
        \qquad
        S=\operatorname{diag}(0.010,0.020,0.010).
    \]
    These make control effort relatively cheap, so the policy can maneuver more
    aggressively.
    \item Strong terminal altitude and velocity penalties:
    \[
        20p_z^2 + 20v_x^2 + 20v_y^2 + 20v_z^2
    \]
    appear in both types' single-player terminal costs.
    \item Type-dependent horizontal anisotropy:
    type $0$ weights $p_y$ heavily and $p_x$ weakly, while type $1$ weights
    $p_x$ heavily and $p_y$ weakly.
    \item Amortization over contexts:
    the learned network is not tied to a single initial state or prior; it is
    trained over random perturbations of both players' initial positions and
    random beliefs.
\end{itemize}

Put together, the experiment learns a context-conditioned public signaling
policy for a 3D landing game in which both players begin above the landing
plane, terminal velocity is explicitly discouraged, and controls are cheap
enough that the learned trajectories can respond strongly to the hidden type,
belief, and initial geometry.

\end{document}
